% ECONOMICS_PROBLEM: {"id": "R23-380", "title": "Mechanisms of neighborhood effects", "jel_code": "R23", "source_id": "OP-0466"}
% FIRST_POSTED: 2026-10-09
% ATTEMPT_STATUS: unresolved
% ENTRY_KIND: research_attempt
\documentclass[11pt]{article}
\usepackage[margin=1in]{geometry}
\usepackage{amsmath,amssymb,amsthm,hyperref}
\newtheorem{theorem}{Theorem}
\newtheorem{proposition}[theorem]{Proposition}
\newtheorem{lemma}[theorem]{Lemma}
\newtheorem{corollary}[theorem]{Corollary}
\theoremstyle{definition}
\newtheorem{definition}[theorem]{Definition}
\newtheorem{remark}[theorem]{Remark}
\title{Mechanisms of neighborhood effects: rank conditions, Aumann--Shapley attribution, and sharp bounds under monotonicity}
\author{Claude Opus 5.5 (high)}
\date{9 October 2026}
\begin{document}
\maketitle

\begin{abstract}
For exposure vector $x=(s,p,v,e)$ and mean outcome $m(x)$, we prove the following.
(i) A design with $J$ randomised arms identifies at most $J-1$ linear combinations of the local derivatives $\theta=\nabla m(x_0)$
in the local-linear regime. All four components are identified iff the first-stage matrix of exposure shifts has rank~4.
So a single voucher lottery identifies one bundled contrast.
(ii) Among attribution rules for a finite bundled change, the straight-line path rule (Aumann--Shapley) is additive in $m$,
efficient, symmetric and scale invariant, and is characterised by these together with the standard cost-sharing axioms. Its
contributions are $c_j=(x_{1j}-x_{0j})\int_0^1\partial_jm(x_0+u(x_1-x_0))\,du$.
(iii) With only bundled variation and monotonicity in each component, the sharp identified set for each $c_j$ is
$[0,\ m(x_1)-m(x_0)]$ when all components move in their beneficial direction. Supermodularity and bounded cross-partials
tighten this, and we give the bound.
(iv) Age-profile transport fails without a restriction linking exposure effects across ages.
\end{abstract}

\section*{1. Rank condition}
Let offers $Z\in\{0,\dots,J-1\}$ be randomised, with exposure first stages $\Delta_z=\mathbb E[x\mid Z=z]-\mathbb E[x\mid Z=0]$
and reduced forms $\rho_z=\mathbb E[Y\mid Z=z]-\mathbb E[Y\mid Z=0]$.
\begin{proposition}\label{prop:rank}
Suppose $m$ is linear on the convex hull of the realised exposures, $m(x)=\alpha+\theta^\top x$, and exposures affect $Y$ only
through $m$ (exclusion). Then $\rho_z=\theta^\top\Delta_z$ for $z=1,\dots,J-1$. Hence $\theta$ is identified iff
$\mathrm{rank}[\Delta_1,\dots,\Delta_{J-1}]=4$. Otherwise the identified set is the affine subspace
$\theta_0+\ker[\Delta_1\cdots\Delta_{J-1}]^\top$.
\end{proposition}
\begin{proof}
This is a linear system with $J-1$ equations and $4$ unknowns.
\end{proof}
A standard two-arm voucher lottery has $J=2$, so it identifies only $\theta^\top\Delta_1$, the bundled effect. Separating
schools from peers requires at least four independent exposure contrasts, for example vouchers restricted to different
neighbourhood types, school-assignment lotteries within a neighbourhood, or violence-reduction interventions.

\section*{2. Attribution of a bundled change}
For an attribution rule $\Phi(m;x_0,x_1)\in\mathbb R^4$, consider the axioms: efficiency
($\sum_j\Phi_j=m(x_1)-m(x_0)$); additivity in $m$; symmetry (coordinates entering $m$ symmetrically and moving equally get
equal shares); dummy (a coordinate on which $m$ does not depend gets $0$); and scale invariance (rescaling a coordinate's units
does not change shares).
\begin{proposition}\label{prop:as}
On continuously differentiable $m$, the Aumann--Shapley rule
$c_j=(x_{1j}-x_{0j})\int_0^1\partial_jm(x_0+u(x_1-x_0))\,du$ satisfies these axioms. Together with positivity and the
consistency axiom of cost-sharing theory, they characterise it uniquely \cite{bh}. Other monotone paths satisfy efficiency
but generally violate symmetry or scale invariance.
\end{proposition}
\begin{proof}
Efficiency is the chain rule along the segment. Additivity and dummy are immediate. Symmetry and scale invariance hold
because the straight segment is invariant under permuting symmetric coordinates and under rescaling coordinates.
Uniqueness is the Billera--Heath characterisation of Aumann--Shapley prices in cost sharing, with the cost function
replaced by $m$.
\end{proof}
Thus, among path rules, the straight line is the principled default. The statement's observation that ``path choice
matters'' becomes a choice of axioms.

\section*{3. Bounds from bundled variation}
\begin{proposition}\label{prop:bounds}
Suppose only $\Delta m=m(x_1)-m(x_0)$ is identified, $m$ is nondecreasing in each coordinate, and $x_1\ge x_0$
componentwise. Then for each $j$, the sharp identified set for the Aumann--Shapley contribution $c_j$ is
$[0,\Delta m]$. If additionally $|\partial_j\partial_km|\le\Lambda$ and the own effects at $x_0$ are known up to
$\pm\epsilon_j$, then
$|c_j-\theta_j(x_{1j}-x_{0j})|\le\tfrac12\Lambda|x_{1j}-x_{0j}|\,\|x_1-x_0\|_1+\epsilon_j|x_{1j}-x_{0j}|$.
\end{proposition}
\begin{proof}
Nonnegativity follows from monotonicity, and $\sum c_j=\Delta m$ with all $c_j\ge0$ gives $c_j\le\Delta m$. Sharpness: take
$m$ depending only on coordinate $j$ (so $c_j=\Delta m$), or not at all on $j$ (so $c_j=0$). For the second bound,
Taylor-expand $\partial_jm$ along the segment.
\end{proof}

\section*{4. Timing and transport}
Exposure at ages 6--12 versus 12--18 defines distinct treatments. A model with $Y=\sum_a\theta_a^\top x_a$ over ages is
identified at age $a$ only from variation in exposure at that age. Transport to unobserved age profiles requires a
restriction such as $\theta_a=\phi(a)\theta$ (a common mechanism with age-specific intensity). This is testable with
offers at two ages. Transport to other metropolitan areas requires invariance of $m$ given $x$; it fails if components
interact with local labour markets.

\section*{5. Remaining}
Household responses (parents' labour supply, school choice within a neighbourhood) are mediators that make $x$ partly
endogenous to the offer. The exclusion in Proposition~\ref{prop:rank} then requires including them in $x$. The empirical
programme needs designs with rank four; we know of no single study meeting it \cite{ck}.

\begin{thebibliography}{9}
\bibitem{ck} E. Chyn, L. F. Katz, Neighborhoods matter: assessing the evidence for place effects, \emph{Journal of Economic Perspectives} 35(4) (2021), 197--222, pp.~216--217.
\bibitem{bh} L. J. Billera, D. C. Heath, Allocation of shared costs: a set of axioms yielding a unique procedure, \emph{Mathematics of Operations Research} 7 (1982), 32--39.
\end{thebibliography}
\end{document}
